\section{Evaluation error budgets}
\label{sec:evaluation}
The theorem in this section is a reusable implication from delivered-data
contracts to field accuracy. Sections~\ref{sec:providers}--\ref{sec:effective}
subsequently construct those data and prove an effective total result.

\subsection{Normalize first and allocate the total error}
Take \(p=\lfloor L/2\rfloor\), \(K=2p\).
Increase the strict-margin onset in \cref{thm:growing} so that
the mathematical truncation uses at most
\(\tfrac12\e^{-L^2/4}\) of the outer budget.
By \eqref{eq:outer-conversion} a sufficient center-normalized
evaluation tolerance is
\begin{equation}
 \varepsilon_X=\frac{c_{\mathrm o}}{2C_\chi}\e^{-3L^2/8}.
 \label{eq:epsilonX}
\end{equation}
Write \(\beta_{\epsilon n}=b_{\epsilon n}/g\) and
\(P_{\epsilon,p}(q)=\sum_{m=0}^{2p}t_{\epsilon m}q^m\).
The normalized exact expression is
\begin{align}
 F/g={}&\sum_{\epsilon,n\le N}\beta_{\epsilon n}
       \e^{\lambda_{\epsilon n}d+T\lambda_{\epsilon n}^2}
       \sum_{m=0}^{2p}t_{\epsilon m}
              \calH_m(d+2T\lambda_{\epsilon n},T)\nonumber\\
 &+\int_{X-V}^{X+V}k_T(z-u)Q_{X,K}(u)/g\dd u.
 \label{eq:normalized-evaluation}
\end{align}
Let \(A_0\) bound its combined center-scaled channel magnitudes.
If \(\abs{\widehat g/g-1}\le\varepsilon_g\le1\), it suffices that
\begin{equation}
 \varepsilon_g A_0\le\varepsilon_X/2,\qquad
 E_{\mathrm{ar}},E_{\mathrm{RS}},E_{\mathrm{quad}},
 E_{\mathrm{geo}},E_{\mathrm{round}}\le\varepsilon_X/20.
 \label{eq:allocation}
\end{equation}
Indeed the center-normalized output error is bounded by
\(\varepsilon_gA_0+(1+\varepsilon_g)\sum E_i\le\varepsilon_X\).
The five categories are, respectively, arithmetic inputs and
preprocessing, retained-integral data, exact quadrature, geometry,
and remaining operation rounding. Scale error is separate.

\subsection{Arithmetic input conditioning}
Use the moments at \(\tau=1,B=2\), and set
\[
 \mathfrak M_m=48M_m+M_m^t,\quad
 \mathcal D_p=\mathsf H_p[\mathfrak M],\quad
 \mathcal D_p^{(1)}=\mathsf H_p[q\mathfrak M],\quad
 \mathcal S_p=\sum_{m=0}^{2p}\mathfrak M_m.
\]
The notation \(\mathsf H_p[q\mathfrak M]\) pairs its \(q^m\)
coefficient with \(\mathfrak M_{m+1}\).
The factor \(48\) includes the buffered spatial Cauchy costs.

\begin{lemma}[Finite coefficient sensitivity]\label{lem:sensitivity}
For \(1\le r\le p\),
\begin{equation}
 \partial_{c_r}P_p(q)=q^{r+1}P_{p-r}(q).
 \label{eq:coefficient-sensitivity}
\end{equation}
For perturbations \(\abs{\delta c_r}\le\varepsilon_c\le1\),
let \(\mathsf C_j^+\) denote \(\mathsf C_j\) with \(A\) replaced
by \(A+1\), and \(\mathsf C_0^+=1\). Then a sufficient condition
sum is
\begin{equation}
 \mathcal C_p=\sum_{r=1}^p\sum_{j=0}^{p-r}
                     \mathsf C_j^+[q^{r+1}\mathfrak M].
 \label{eq:condition-sum}
\end{equation}
\end{lemma}
\begin{proof}
Differentiate the local generating function: differentiation
in \(c_r\) multiplies it by \(q^{r+1}s^r\), proving
\eqref{eq:coefficient-sensitivity} coefficient by coefficient.
Along the segment of perturbed coefficients,
\(\abs{c_r'}\le(A+1)D^{r-1}\), after dropping the favorable
\(X^{-r}\). Apply the positive partition majorant and integrate
the derivative along that segment. Then use the Gaussian
channel moments, giving \eqref{eq:condition-sum}.
\end{proof}

Impose
\(D_0/X\le1/2-\pi/8\) and \(\varepsilon_\lambda\le1/2\).
The sectorial slope formula ensures that the original slopes
and their perturbed connecting segments have real part of
absolute value at most one.
For
\[
 \abs{\delta\beta_{\epsilon n}}\le\varepsilon_\beta n^{-1/2},
 \quad\abs{\delta\lambda_{\epsilon n}}\le\varepsilon_\lambda,
 \quad\abs{\delta c_{\epsilon r}}\le\varepsilon_c,\quad
 \abs{\delta t_{\epsilon m}^{\mathrm{store}}}\le\varepsilon_t,
\]
one has
\begin{equation}
\begin{split}
 E_{\mathrm{ar}}\le C_{\mathrm{pop}}X^{1/4}
 \bigl[\varepsilon_\beta\mathcal D_p
 +(1+\varepsilon_\beta)
 \{\varepsilon_\lambda\mathcal D_p^{(1)}
               +\varepsilon_c\mathcal C_p+\varepsilon_t\mathcal S_p\}\bigr].
\end{split}
\label{eq:arithmetic-error}
\end{equation}
To verify the estimate, change \(\beta\) first with the exact
polynomial and slope, then the slope, then the shape coefficients
and stored output. Differentiation in \(\lambda\) inside the
real Gaussian transform inserts \(q\).
The real tilt remains bounded by \(\e^{\abs q}\), so the same
moments apply. This ordering avoids using an unperturbed
polynomial majorant for perturbed coefficients.
The error \(\varepsilon_t\) includes the certified error of
preprocessing at the supplied \(c_r\); it is charged here
only once.

The logarithms of all finite condition sums are
\(O((p+1)\log(p+2))\), by the factorial bounds and the finite
partition formulas. There is no division by a possibly vanishing
field or correction.

\subsection{Large phases and very small scales}
At real \(X\),
\[
 \beta_{\epsilon n}=n^{-1/2}
    \e^{\ii(\theta_\epsilon(X)-\epsilon X\log n/2)},\qquad
 g=\e^{\Rea\ell_+(X)}.
\]
If the error in the carrier logarithm is at most \(q_{\mathrm{err}}\),
then
\(\abs{\delta\beta_{\epsilon n}}
\le n^{-1/2}\e^{q_{\mathrm{err}}}q_{\mathrm{err}}\).
A sufficient allowance is
\begin{align}
 q_{\mathrm{err}}\le{}&
 \abs{\delta\theta}
 +(X/2+1/2)\abs{\delta\log n}
 +\tfrac12\log n\abs{\delta X}
 +\tfrac12\abs{\delta X\,\delta\log n}\nonumber\\
 &+\abs{\delta(\Rea\log M-\log g)}.
 \label{eq:phase-budget}
\end{align}
Also
\(\abs{\delta\lambda}\le\abs{\delta\ell'}+
\abs{\delta\log n}/2\) at fixed \(X\).
Thus logarithms of mode indices must be accurate to order
\(\varepsilon_\beta/X\).
Unreduced phases have size \(O(XL)\), requiring
\(O(L+\log L)\) guard bits when formed by relative arithmetic.
Normalized gamma ratios should be formed from log differences.
Dividing separately underflowed values of \(M\) and \(g\)
does not implement this specification.
Sufficient exponent range is part of the model.

\subsection{Real-node data and geometry}
At a real node on branch \(m\), write
\[
 A=(-1)^{m-1}M(w_+(u))U(u)a(u)^{-1/2}/g,\qquad R_k=a(u)^{-k}.
\]
If the errors in \(A,R_k,C_k\), and the normalized transfer
are bounded by \(\varepsilon_A,\varepsilon_{R_k},
\varepsilon_{C_k},\varepsilon_{\mathrm{tr}}\), respectively, then
\begin{equation}
\begin{split}
 \abs{\widehat Q_m/g-Q_m/g}
 \le\varepsilon_{\mathrm{tr}}+2\sum_{k=0}^K
 \bigl[&\varepsilon_A R_k\abs{C_k}\\
 &+(\abs A+\varepsilon_A)
 \{\varepsilon_{R_k}\abs{C_k}
           +(R_k+\varepsilon_{R_k})\varepsilon_{C_k}\}\bigr].
\end{split}
\label{eq:node-data-error}
\end{equation}
This is just the product rule for finite perturbations.
It retains every coefficient and the exact transfer.

Geometry errors use a local strip bound rather than the much
larger branch-ellipse bound. Require \(V\ge3\),
\[
 4(V+2)/\sqrt{2\pi X}\le1/2,\qquad 2(V+2)\le X/4.
\]
Define
\begin{subequations}\label{eq:local-strip}
\begin{align}
 H_M^{\mathrm{loc}}
 &=\exp\left\{\frac{\pi}{8}(V+2)+\frac L4+
             \frac{D_0(V+3)}X+\frac{H(V+3)^2}{2X}\right\},\\
 H_U^{\mathrm{loc}}&=\e^{(L+1)/4},\\
 B_Q^{\mathrm{loc}}
 &=2H_M^{\mathrm{loc}}\left[
       \sum_{n\in\mathcal T} n^{-1/2}n^{1/2}
       +H_U^{\mathrm{loc}}a_*^{-1/2}\sum_{k=0}^KD_ka_*^{-k}\right],\\
 K_{\mathrm{loc}}
 &=\frac{\e^{(Y+1)^2/4}}{\sqrt{4\pi}}
       \left[\max_{a\le3}\mathsf K_a(W+V+4)
                         +\mathsf K_2(W+V+4)\right],\\
 B_{\mathrm{loc}}&=B_Q^{\mathrm{loc}}K_{\mathrm{loc}}.
\end{align}
\end{subequations}
On \(\abs{\Rea u-X}\le V+2,\abs{\Ima u}\le1\), these bound
the normalized integrand channels. For gamma, use
\(\ell_+'(X)=-\pi/8+\ii L/4+O(D_0/X)\) and bound real and
imaginary displacements separately. This proves the displayed
\(H_M^{\mathrm{loc}}\); using \(\abs{\ell'}\) times the full
window instead would lose the useful bound
\[
 \log^+B_{\mathrm{loc}}=O(L)+O(K\log(K+2)).
 \label{eq:local-strip-growth}
\]
Radius-\(1/2\) Cauchy estimates bound each derivative in \(u\)
by twice its local bound.

Let \(\varepsilon_w\) be the sum of absolute physical weight
errors, and let real approximate nodes have maximum displacement
\(\varepsilon_u\le1/4\). Use their fixed reference branch labels.
If omitted boundary strips have total length at most
\(6\varepsilon_b\), \(0\le\varepsilon_b\le1/4\), then
\begin{equation}
 E_{\mathrm{geo}}\le B_{\mathrm{loc}}
   \{\varepsilon_w+2(2V+\varepsilon_w)\varepsilon_u
                                      +6\varepsilon_b\}.
 \label{eq:geometry-error}
\end{equation}
The terms follow by changing weights, then nodes, and integrating
over omitted strips. Such strips must lie inside the true fixed
branches. There are at most three boundaries, so two-sided
width-\(\varepsilon_b\) strips have the stated total length.
For exact reference endpoints take \(\varepsilon_b=0\);
endpoint errors propagated into node and weight errors are not
charged again. Rules extending outside the original window need
a new length and ellipse analysis.

If \eqref{eq:node-data-error} is bounded by \(\varepsilon_Q\)
at every approximate node, then the separate data cost is
\begin{equation}
 E_{\mathrm{RS}}\le(2V+\varepsilon_w)K_{\mathrm{loc}}\varepsilon_Q.
 \label{eq:RS-data-budget}
\end{equation}
An explicit scale bound for \eqref{eq:allocation} is
\begin{equation}
 A_0=C_{\mathrm{pop}}X^{1/4}\mathsf H_p[\mathfrak M]
                                      +2VB_{\mathrm{loc}}.
 \label{eq:A0}
\end{equation}
Its logarithm is \(O(L\log L)\).

\subsection{A coefficient construction away from apparent poles}
Let \(h_m(\chi)=F_0^{(m)}(\chi)/m!\).
On the rectangle \(\Gamma\) with vertices \(\pm2\pm\ii\),
Cauchy's formula gives
\begin{equation}
 h_m(\chi)=\frac1{2\pi\ii}\oint_\Gamma
                         \frac{F_0(\zeta)}{(\zeta-\chi)^{m+1}}\dd\zeta.
 \label{eq:contour-jet}
\end{equation}
For \(\abs{\Rea\chi}\le3/2,\abs{\Ima\chi}\le1/2\), use
a Gauss rule with \(q_c\) nodes on each straight edge and ellipse
parameter \(\rho_c=9/8\).
The edge ellipses have semiaxis ratios \(145/144,17/144\);
they satisfy
\[
 \abs{\Rea\zeta}\le9/4,\quad\abs{\Ima\zeta}\le5/4,\quad
 \abs{2\cos(\pi\zeta)}\ge1,\quad
 \operatorname{dist}(\zeta,\chi)\ge1/4.
\]
For the vertical edges the real deviation from \(\pm2\) is
at most \(17/144\); for the horizontal edges the imaginary
distance from the real line is at least \(1-34/144\).
These give the denominator bounds and the stated separation.
With \(B_c=\e^{45\pi/16}+\sqrt2\cosh(5\pi/8)\), the integrand
is bounded by \(B_c4^{m+1}\).
The total contour length is \(12\). Applying \cref{lem:gauss}
and the factor \(1/(2\pi)\) proves
\begin{equation}
 \varepsilon_{h,m}^{\mathrm{quad}}
 \le\frac{24}{\pi}B_c4^{m+1}
                  \frac{\rho_c^{-2q_c+1}}{\rho_c-1}.
 \label{eq:contour-error}
\end{equation}
For all \(m\le3K\), a requested error
\(\exp[-3L^2/8-O(L\log L)]\) is reached with
\(q_c=O(L^2+L\log L)\).
The identities now take the normalized form
\[
 C_k=\pi^{-2k}\sum_j(\pi/(2\ii))^j e_j^{(k)}h_{3k-2j}.
\]
Each perturbation term is bounded by
\(\abs{e_j}\varepsilon_h+\abs{h_{3k-2j}}\varepsilon_e+
\varepsilon_e\varepsilon_h\), times the displayed powers of \(\pi\).
Errors in those powers also require their product bounds.
This construction supplies an analytical way around the apparent
poles; primitive, node, weight, and rounding errors still have to
be certified. The fixed contour samples can be reused for all
outer quadrature nodes.

\subsection{Rounding and sufficient requested precision}
Use real-component arithmetic with unit roundoff \(u\), sufficient
exponent range, and certified bounds \(\abs a\le A_a\),
\(\abs b\le A_b\), with propagated errors \(E_a,E_b\).
For addition, multiplication, and division by an exact positive
integer \(j\), valid updates are
\begin{subequations}\label{eq:roundoff-updates}
\begin{align}
 E_+&=E_a+E_b+u(A_a+A_b+E_a+E_b),\\
 E_\times&=A_aE_b+A_bE_a+E_aE_b
                          +u(A_a+E_a)(A_b+E_b),\\
 E_{\div j}&=[E_a+u(A_a+E_a)]/j.
\end{align}
\end{subequations}
Noninteger reciprocals require certified separation and their
own error. Nonlinear primitives are not automatically covered
by these arithmetic updates. For example, if real exponential
relative error and sine/cosine absolute error are at most \(u\),
then, for \(u\le1/16\), a complex exponential with exact real
part at most \(C_0\) and argument error \(E_{\mathrm{arg}}\) has
error at most
\[
 \e^{C_0}(\e^{E_{\mathrm{arg}}}-1)
                  +4u\e^{C_0+E_{\mathrm{arg}}}.
\]

The exact derivative channels are also finite recurrences:
\begin{subequations}\label{eq:channel-formulas}
\begin{align}
 \partial_d^a(\beta\e^{\lambda d+T\lambda^2}\calH_m)
 &=\beta\e^{\lambda d+T\lambda^2}
   \sum_{j\le\min(a,m)}\binom aj(m)_j
                  \lambda^{a-j}\calH_{m-j},\\
 \partial_T(\beta\e^{\lambda d+T\lambda^2}\calH_m)
 &=\beta\e^{\lambda d+T\lambda^2}
   [\lambda^2\calH_m+2m\lambda\calH_{m-1}
                           +m(m-1)\calH_{m-2}].
\end{align}
\end{subequations}
All \(\calH\)'s here have argument \(d+2T\lambda,T\), and
\((m)_j=m!/(m-j)!\).
Evaluation needs retained moments only through \(2p\).
The error analysis still needs moments through \(2p+4\).

To bound intermediate magnitudes, take
\(\Lambda=1+L/4\ge\abs\lambda\) after an additional slope gate,
\(t_*=W+4\Lambda\), and
\[
 h_0^*=1,\quad h_1^*=t_*,\quad
 h_m^*=t_*h_{m-1}^*+4(m-1)h_{m-2}^*.
\]
Completing the square in \(\Ima\lambda\), with
\(\abs{\Rea\lambda}\le1\), gives
\[
 \abs{\e^{\lambda d+T\lambda^2}}
                  \le\e^{R+2+Y^2/4}.
\]
Thus no artificial \(\e^{+L^2}\) carrier bound is required
in the actual real-node circuits.
The triangular coefficient recurrence, moment recurrence,
weighted Riemann--Siegel recurrence, and finite linear sums
have positive path majorants of size
\((N+1)\exp(O(L\log L))\) times a polynomial in the node counts.
One can see this by expanding the triangular recurrences:
their depth is \(O(p+K)\), their factors are polynomial in
\(p,K,L\), and their partition paths have
\(\exp(O((p+K)\log(p+K+2)))\) total weight.
Linear accumulation contributes its actual length, not
an exponential in \(N\). Impose
\(uS_{\mathrm{lin}}\le1/4\), where \(S_{\mathrm{lin}}\)
is the longest actual linear accumulation count.

\begin{theorem}[Modular total evaluation accuracy]
\label{thm:evaluation}
Assume the half-target truncation onset, the displayed slope and
geometry gates, exact query coordinates, and exact frozen
descriptors or a certified guard-strip procedure. Suppose
certified inputs, nodes, weights, scale, and the arithmetic
recurrences satisfy
\cref{eq:allocation,eq:arithmetic-error,eq:node-data-error,eq:geometry-error,eq:RS-data-budget,eq:roundoff-updates}.
Then the separately computed field, spatial-through-three, and
first-time outputs have total outer-envelope error at most
\(\e^{-L^2/4}\).
A sufficient asymptotic requested normalized accuracy is
\begin{equation}
 P_{\mathrm{req}}=
             \frac{3L^2}{8\log2}+O_{r,b}(L\log L)
 \quad\hbox{bits}.
 \label{eq:precision}
\end{equation}
The finite displayed inequalities, rather than the \(O\)-term
alone, specify acceptance at a chosen height.
\end{theorem}
\begin{proof}
The truncation uses half the target.
The allocation proof after \eqref{eq:allocation}, followed by
\eqref{eq:outer-conversion}, bounds the evaluation error by
the other half.
For \eqref{eq:precision}, every real-node conditioning and
intermediate-size cost has logarithm \(O(L\log L)\).
Also \(\log N=L/2+O(1)\) and both quadrature orders are
\(O(L^2)\). The positive error recurrences therefore require
normalized input errors and roundoff can be made sufficiently
small by choosing them at most
\(\exp[-3L^2/8-C_{r,b}L\log L]\), with \(C_{r,b}\)
large enough for their finite majorants.
Taking base-two logarithms gives \eqref{eq:precision}.
The phase guards in \eqref{eq:phase-budget} and the small
accumulation guard are lower-order requirements within that
choice.
\end{proof}
The precision is sufficient, not necessary. It specifies
delivered input accuracy and the declared arithmetic model;
it does not bound a concrete provider's internal precision,
iteration count, or bit complexity. Those costs are supplied separately
in \cref{sec:providers,subsec:bit-complexity}.
For arbitrary real coordinates, exact floor decisions can require
arbitrarily fine discrimination near a threshold. Rounded
floor tests are not part of the theorem. Uncertainty in \(z,T\)
also needs its own derivative-based error allocation.

\subsection{Costs and the meaning of an exact residual}
With named data already supplied, polynomial preprocessing uses
\(O(p^3)\) scalar arithmetic, and all arithmetic channels use
\(O(Np)\). The weighted Riemann--Siegel table uses \(O(K^2)\).
The contour construction at all outer nodes uses
\(O(n(q_cK+K^2))=O(L^5)\) scalar arithmetic, with fixed contour
values reused. This last count is for the completion part only;
it does not replace the \(N\asymp\sqrt X\) arithmetic modes.
All counts exclude provider bit costs.

Three residual statements must remain distinct.
The exact frozen model solves heat exactly.
A symbolic model with fixed perturbed coefficients and fixed
quadrature nodes, weights, and node data also solves heat
exactly, since it is a sum of heat solutions.
Independently rounded channel outputs have only a bounded
defect: their \(T\)-minus-\(zz\) discrepancy is bounded by
the sum of the paid channel errors.
A rounded spatial derivative of that defect would require
an additional mixed \(T,z\) output channel. It is not supplied
by differentiating a pointwise rounding allowance.
